\section{Dataset Construction Details}
\label{app:data}

This appendix documents the data actually used in the pilot experiments
(\S\ref{sec:results}); steps that belong to the final protocol but are not yet
exercised by the pilot are explicitly marked \emph{planned}.

\subsection{Source and input construction (pilot)}
All pilot examples come from BigVul \citep{fan2020bigvul}, which provides, for
each vulnerability-fixing commit in C/C++ projects, the vulnerable function
(\texttt{func\_before}), the fixed function (\texttt{func\_after}), the commit
message, and a CVE-derived CWE label. We keep rows marked vulnerable whose CWE
normalizes to a leaf of our tree (below), and render each example as a
\emph{unified diff} from the vulnerable to the fixed function, prefixed by an
instruction comment (``\texttt{'-'} lines are the VULNERABLE code \ldots
classify the weakness in the removed/vulnerable code'') and the first 500
characters of the commit message. The diff localizes the vulnerable lines and is
the input representation under which models clear the $\approx\!40\%$ category
ceiling we observed on raw functions. Diffs with fewer than three lines are
dropped; code is truncated at 5{,}000 characters at prompt time
(\S\ref{app:prompts}).

\subsection{Splits (pilot)}
The dev set ($n{=}50$) is drawn from the BigVul validation split and the pilot
set ($n{=}500$) from the BigVul test split --- so prompt iteration on dev cannot
overfit the pilot --- using per-CWE stratified round-robin sampling (seed 42).
Pilot statistics: 500 examples, 18 distinct CWE leaves, 73 distinct projects;
median diff length 15 lines (mean 27). Category composition:
\textsc{buffer} 214, \textsc{injection} 117, \textsc{access} 77,
\textsc{lifetime} 54, \textsc{numeric} 28, \textsc{resource} 5,
\textsc{control} 5; the largest leaves are CWE-119 (128), CWE-20 (112),
CWE-125 (68), CWE-200 (49). Most reported runs use the first 300 (nano) or 150
(Haiku, DeepSeek) examples of this fixed, shuffled file; per-example outputs are
released for paired analyses.

\subsection{CWE trees}
\paragraph{Full tree.} We download the official MITRE CWE corpus
(\texttt{cwec\_latest.xml}), extract the Research-Concepts view (View-1000)
ChildOf hierarchy rooted at the ten Pillars, and resolve multi-parent nodes to a
\emph{primary} parent by breadth-first search from the pillars (first, i.e.\
minimal-depth, discovery wins). The result has 944 nodes and 716 leaves (depths
0--5) and is released as adjacency JSON with SHA-256
\texttt{7aeff209\ldots{}26948a6}.

\paragraph{Coarse code-aligned tree (pilot).} MITRE pillars are research
abstractions (``Improper Control of a Resource Through its Lifetime'') that even
strong policies place poorly, which floors level-0 accuracy and makes the
policy-strength axis (RQ2) unreadable. For the pilot we therefore regroup the
22 dataset CWEs into a 2-level tree with 7 code-recognizable super-categories
(spatial memory, temporal memory / lifetime, resource exhaustion, numeric,
injection / input handling, control-flow, access control / info exposure /
misc); leaves keep their MITRE names and descriptions. All $26$ leaves are
members of CWE View-1003, the simplified-mapping view that NVD uses to label
published CVEs, so this label space is a subset of the official operational
vocabulary rather than an ad-hoc simplification. The tree is released with
SHA-256 \texttt{f3013c19\ldots{}9c79aa}. The full-depth tree is used only for
the granularity boundary-condition probe (\S\ref{sec:granularity}).

\subsection{Multi-source pooling and deduplication (planned)}
The final protocol pools PrimeVul \citep{ding2024primevul}, DiverseVul
\citep{chen2023diversevul}, BigVul, and CVEfixes \citep{bhandari2021cvefixes},
with function-level MinHash deduplication (Jaccard $\ge 0.85$) and
repository-level disjointness across splits; the pilot uses a single source
(BigVul) with splits inherited from the dataset's own validation/test
partition, so cross-source leakage is not yet a concern at pilot scale.

\subsection{Pretraining-leakage slice (planned)}
Because public CVEs may appear in pretraining corpora, the final protocol
reports a slice restricted to CVEs published after each model's training
cutoff. At pilot scale we note that all four policies see the same examples, so
leakage would shift absolute numbers rather than the strategy comparisons that
are the object of study.

\subsection{Datasheet (planned)}
A Datasheets-for-Datasets \citep{gebru2021datasheets} entry (motivation,
composition, collection, preprocessing, uses, distribution, maintenance) will
accompany the released benchmark bundle.
